%**************************************************************************************
% บทคัดย่อภาษาไทย
%**************************************************************************************
\phantomsection
\addcontentsline{toc}{chapter}{บทคัดย่อภาษาไทย}

\vspace*{-0.85in}

{\renewcommand{\arraystretch}{0.75}
\begin{tabularx}{\textwidth}{l L}
ชื่อ                           & : นายวีรภัทร ชื่นชอบ\\
                             & : นายธนากร ธนบวรพล\\
ชื่อปริญญานิพนธ์                  & : ปัญญาประดิษฐ์แชทบอทสำหรับองค์กร กรณีศึกษา ฮาร์ทวอมมิ่ง จำกัด\\
สาขาวิชา                       & : เทคโนโลยีสารสนเทศ\\
สถาบัน                         & : มหาวิทยาลัยเทคโนโลยีพระจอมเกล้าพระนครเหนือ\\
อาจารย์ที่ปรึกษาปริญญานิพนธ์        & : อาจารย์ ดร.ศิรินทรา แว่วศรี\\
ปีการศึกษา                     & : 2569\\
\end{tabularx}
}

\vspace{-2mm}
\begin{center}\textbf{บทคัดย่อ}\end{center}
\vspace{-3mm}

\hspace*{1.5em}ปริญญานิพนธ์นี้มีวัตถุประสงค์เพื่อวิเคราะห์ ออกแบบ และพัฒนาระบบปัญญาประดิษฐ์แชทบอทสำหรับองค์กร กรณีศึกษา บริษัท ฮาร์ทวอมมิ่ง จำกัด ผู้นำเข้าและจัดจำหน่ายสินค้าไลฟ์สไตล์พรีเมียม โดยนำเทคโนโลยีปัญญาประดิษฐ์มาแก้ไขปัญหาข้อมูลกระจัดกระจาย ความล่าช้าในการสืบค้นข้อมูลสินค้าและระเบียบปฏิบัติ ความผิดพลาดในการบันทึกบัญชีรายรับ-รายจ่ายจากสลิปโอนเงินและใบเสร็จ ตลอดจนความยุ่งยากในการบริหารงานบุคคล เช่น การยื่นใบลาและการอนุมัติงาน

\hspace*{1.5em}ระบบที่พัฒนาขึ้นผสานการทำงานของ LINE Official Account ผ่าน LINE Messaging API ร่วมกับเว็บแอปพลิเคชันส่วนกลาง พัฒนาด้วย Python Flask, SQLite และ ChromaDB โดยในภาคการศึกษาที่ 1/2569 ได้ดำเนินการวิเคราะห์ ออกแบบระบบ และพัฒนาระบบส่วนโครงสร้างพื้นฐาน ได้แก่ 1) ระบบยืนยันตัวตนและจัดการบัญชีผู้ใช้งาน (Email/Password, Google OAuth 2.0, LINE Linking) 2) ระบบคลังความรู้อัจฉริยะ (RAG) เชื่อมต่อโมเดลภาษาขนาดใหญ่ Google Gemini เพื่อสืบค้นและตอบคำถามภาษาธรรมชาติจากเอกสารองค์กร (PDF, TXT, สเปรดชีต) 3) ระบบค้นหาเวกเตอร์ (Vector Search) ด้วย ChromaDB และ FastEmbed และ 4) ระบบจัดการเอกสารเชื่อมโยง Google Drive จัดระเบียบไฟล์แยกโฟลเดอร์อัตโนมัติ

\hspace*{1.5em}ผลการดำเนินงานในภาคการศึกษาที่ 1/2569 พบว่า ระบบโครงสร้างพื้นฐานและโมเดลคลังความรู้ที่พัฒนาขึ้นสามารถทำงานได้ครบถ้วนตามขอบเขตที่กำหนด โดยผ่านการทดสอบการทำงานเบื้องต้น (Black-Box Testing) ระบบสามารถรับข้อความคำถามผ่าน LINE Official Account และประมวลผลดึงคำตอบจากเอกสารของบริษัทมาตอบกลับในรูปแบบ Flex Message ได้อย่างถูกต้องและตรงตามบริบทของสินค้าจริง สำหรับระบบจัดการการเงินและบัญชีอัตโนมัติด้วยเทคโนโลยี OCR, ระบบบริหารงานบุคคล, แผงควบคุมสถิติภาพรวม และการประเมินความพึงพอใจของผู้ใช้งานจริง จะดำเนินการพัฒนาและประเมินผลในภาคการศึกษาที่ 2/2569 ต่อไป

\vspace{1mm}
\begin{flushright}
(ปริญญานิพนธ์มีจำนวนทั้งสิ้น \pageref{LastPage} หน้า)

\vspace{2mm}
\rule{8 cm}{0.4pt} อาจารย์ที่ปรึกษาปริญญานิพนธ์\\
(อาจารย์ ดร.ศิรินทรา แว่วศรี)
\end{flushright}


%**************************************************************************************
% บทคัดย่อภาษาอังกฤษ
%**************************************************************************************
\newpage
\phantomsection
\addcontentsline{toc}{chapter}{บทคัดย่อภาษาอังกฤษ}

\vspace*{-0.85in}

{\renewcommand{\arraystretch}{0.75}
\begin{tabularx}{\textwidth}{l L}
Name                    & : MR. WEERAPHAT CHUENCHOB\\
                        & : MR. THANAKORN THANABAWONPON\\
Project Title           & : Artificial Intelligence Chatbot for Organization (Case Study: Heartwarming Co., Ltd.)\\
Major Field             & : Information Technology\\
Institution             & : King Mongkut’s University of Technology North Bangkok\\
Project Advisor         & : Dr. Sirintra Vaiwsri\\
Academic Year           & : 2026\\
\end{tabularx}
}

\vspace{-2mm}
\begin{center}\textbf{Abstract}\end{center}
\vspace{-3mm}

\hspace*{1.5em}The objective of this special project is to analyze, design, and develop an Artificial Intelligence Chatbot for Organization, with Heartwarming Company Limited as a case study. The project applies artificial intelligence technologies to address problems regarding scattered enterprise information, delayed product inquiries, manual bookkeeping errors from transaction slips and invoices, and cumbersome human resource management processes such as leave requests and approvals.

\hspace*{1.5em}The system integrates LINE Official Account via LINE Messaging API with a centralized web dashboard, engineered using Python Flask, SQLite, and ChromaDB. In the first semester of academic year 2026 (Semester 1/2026), the project focuses on infrastructure analysis, system design, and core module implementation, comprising: 1) user management and authentication (Email/Password, Google OAuth 2.0, LINE Linking), 2) Retrieval-Augmented Generation (RAG) knowledge base leveraging Google Gemini to answer natural language questions from company documents (PDF, TXT, Spreadsheets), 3) semantic vector search utilizing ChromaDB and FastEmbed, and 4) document management integrated with Google Drive API for automated folder organization.

\hspace*{1.5em}The implementation results for Semester 1/2026 demonstrate that all planned core infrastructure modules operate successfully according to defined specifications, verified through baseline Black-Box testing. The chatbot seamlessly handles incoming queries via LINE OA and synthesizes accurate, context-aware responses presented via interactive LINE Flex Messages. The advanced automated financial bookkeeping with OCR, human resource leave management, executive analytical dashboard, and formal end-user satisfaction evaluation will be subsequently implemented and evaluated in Semester 2/2026.

\vspace{1mm}
\begin{flushright}
(Total \pageref{LastPage} pages)

\vspace{2mm}
\rule{8 cm}{0.4pt} Project Advisor\\
(Dr. Sirintra Vaiwsri)
\end{flushright}


%**************************************************************************************
% กิตติกรรมประกาศ
%**************************************************************************************
\newpage
\phantomsection
\addcontentsline{toc}{chapter}{กิตติกรรมประกาศ}

\vspace*{-0.5in}

\begin{center}\textbf{กิตติกรรมประกาศ}\end{center}
\vspace*{0.2in}

\hspace*{1.5em}ปริญญานิพนธ์ เรื่อง "ปัญญาประดิษฐ์แชทบอทสำหรับองค์กร กรณีศึกษา ฮาร์ทวอมมิ่ง จำกัด" ฉบับนี้ สำเร็จลุล่วงไปได้ด้วยความกรุณาและความช่วยเหลืออย่างดียิ่งจาก อาจารย์ ดร.ศิรินทรา แว่วศรี อาจารย์ที่ปรึกษาปริญญานิพนธ์ ที่ได้กรุณาสละเวลาให้คำปรึกษา แนะนำแนวทางการวิเคราะห์ ออกแบบระบบ และการแก้ไขปัญหาทางเทคนิคต่าง ๆ ตลอดจนตรวจสอบและปรับปรุงรายงานฉบับนี้ให้มีความถูกต้องสมบูรณ์ คณะผู้จัดทำขอกราบขอบพระคุณเป็นอย่างสูงไว้ ณ โอกาสนี้

\vspace{0.8em}
\hspace*{1.5em}ขอขอบพระคุณ ผู้ช่วยศาสตราจารย์ ดร.สุปิติ กุลจันทร์ ประธานกรรมการสอบปริญญานิพนธ์ และผู้ช่วยศาสตราจารย์ ดร.พาฝัน ดวงไพศาล กรรมการสอบปริญญานิพนธ์ ที่ได้ให้ข้อเสนอแนะเชิงวิชาการอันทรงคุณค่า เพื่อให้โครงงานมีความสมบูรณ์รอบด้านยิ่งขึ้น

\vspace{0.8em}
\hspace*{1.5em}ขอขอบพระคุณ ผู้บริหารและพนักงานของ บริษัท ฮาร์ทวอมมิ่ง จำกัด ทุกท่าน ที่ได้ให้ความอนุเคราะห์ข้อมูลจริงในการศึกษา ให้ความร่วมมือในการทดสอบระบบ และให้ข้อเสนอแนะจากการใช้งานจริง ซึ่งเป็นประโยชน์อย่างยิ่งต่อการพัฒนาระบบให้ตอบสนองความต้องการของธุรกิจได้อย่างแท้จริง

\vspace{0.8em}
\hspace*{1.5em}ท้ายที่สุดนี้ คณะผู้จัดทำขอกราบขอบพระคุณ บิดา มารดา และครอบครัว ที่ได้ให้การอุปการะสนับสนุนการศึกษา และเป็นกำลังใจสำคัญเสมอมา ตลอดจนเพื่อน ๆ นักศึกษาสาขาวิชาเทคโนโลยีสารสนเทศ ที่ได้แลกเปลี่ยนความรู้และให้ความช่วยเหลือตลอดระยะเวลาการศึกษา

\vspace{1.5cm}
\begin{flushright}
นายวีรภัทร ชื่นชอบ\\
นายธนากร ธนบวรพล
\end{flushright}
